% !TEX root = main.tex
% ===============================================================
%  zhbooksetting.tex —— zhbook 中文书籍模板 全局设置文件
%  由 main.tex 自动载入 (\input); 标题/作者等封面信息仍在 main.tex 中设置.
%  所有字体均取自本仓库 font/ 目录, 在任何电脑上编译都无需安装字体.
% ===============================================================
% ===============================================================
%  一、字体设置
% ===============================================================

%----- 中文字体: Fandol 字族 (取自 font/) -----
% 与 ctex 官方 fandol 字体集使用同一套字体, 此处改为从本地目录读取文件;
% 字体族名与官方保持一致 (zhsong/zhhei/zhfs/zhkai), 中文标点处理不变.
\setCJKmainfont{FandolSong-Regular.otf}[        % 正文: 宋体
    Path = font/,
    BoldFont = FandolSong-Bold.otf,             % 粗体: 宋体粗
    ItalicFont = FandolKai-Regular.otf]         % 斜体: 楷体
\setCJKsansfont{FandolHei-Regular.otf}[         % 无衬线: 黑体
    Path = font/,
    BoldFont = FandolHei-Bold.otf]
\setCJKmonofont{FandolFang-Regular.otf}[        % 等宽: 仿宋
    Path = font/]
% 字体族声明: 使 \songti \heiti \kaishu \fangsong 均使用本地字体文件
\setCJKfamilyfont{zhsong}{FandolSong-Regular.otf}[Path = font/, BoldFont = FandolSong-Bold.otf]
\setCJKfamilyfont{zhhei}{FandolHei-Regular.otf}[Path = font/, BoldFont = FandolHei-Bold.otf]
\setCJKfamilyfont{zhfs}{FandolFang-Regular.otf}[Path = font/]
\setCJKfamilyfont{zhkai}{FandolKai-Regular.otf}[Path = font/]

%----- 封面字体: 思源黑体 (展示用, 比 Fandol 黑体更现代; 供文类封面使用) -----
\setCJKfamilyfont{zhcover}{SourceHanSansCN-Regular.otf}[Path = font/, BoldFont = SourceHanSansCN-Bold.otf]
\renewcommand{\coverfont}{\CJKfamily{zhcover}}   % 封面标题/副标题/作者均使用该字体

%----- 西文字体 -----
% 衬线 (正文): STIX Two Text —— Times 字体现代修订版, 学术出版风格, 与数学字体同族
\setmainfont{STIXTwoText-Regular.ttf}[
    Path = font/,
    BoldFont = STIXTwoText-Bold.ttf,
    ItalicFont = STIXTwoText-Italic.ttf,
    BoldItalicFont = STIXTwoText-BoldItalic.ttf]
% 无衬线: 思源黑体 (其西文字形为 Source Sans, 干净现代)
\setsansfont{SourceHanSansCN-Regular.otf}[
    Path = font/, BoldFont = SourceHanSansCN-Bold.otf]
% 等宽 (代码/命令行): JetBrains Mono
\setmonofont{JetBrainsMono-Regular.ttf}[
    Path = font/,
    BoldFont = JetBrainsMono-Bold.ttf,
    ItalicFont = JetBrainsMono-Italic.ttf,
    BoldItalicFont = JetBrainsMono-BoldItalic.ttf]

%----- 数学字体: STIX Two Math (取自 font/ 目录, 经 unicode-math 接入) -----
% STIX Two 没有粗体数学字形, 故指定同源字体 XITS Math Bold 充当粗体
% (供向量、矩阵等 \boldsymbol 使用; 如不需要粗体数学, 可去掉 BoldFont 选项)
\RequirePackage{unicode-math}                   % OpenType 数学排版引擎
\setmathfont{STIXTwoMath-Regular.ttf}[
    Path = font/, BoldFont = XITSMath-Bold.otf]
% 证毕符号: 改用数学字体自带的白色方格, 避免使用会回退到 AMS 字体的 \square
\renewcommand{\qedsymbol}{\ensuremath{\mdlgwhtsquare}}

% 如需更换字体: 把字体文件放入 font/ 目录 (或其他自建子目录), 再修改上面相应设置即可.

% ===============================================================
%  二、常用宏包
% ===============================================================
% 代码排版 (listings) 已由文类载入并配置样式 (行号/边框/配色), 可直接使用 lstlisting 环境
\usepackage{subcaption}  % 子图排版 (caption 宏包的配套子图宏包)
%\usepackage[notcite,notref]{showkeys}  % 调试用: 页边显示引用/标签键名

%----- 取消链接颜色和方框 (hyperref 在导言末尾载入, 设置须放入 \AtEndPreamble) -----
%\AtEndPreamble{\hypersetup{hidelinks}}

% ===============================================================
%  三、参考文献 / 引用 / 索引
% ===============================================================
%----- 参考文献格式 -----
\bibliographystyle{gbt7714-2025-numeric}     % GB/T 7714-2025 顺序编码制
%\bibliographystyle{gbt7714-2025-authoryear} % GB/T 7714-2025 著者-出版年制
%\bibliographystyle{plain}  % 其他可选: abbrv, unsrt, siam
%\bibliographystyle{thuthesis-numeric}  % 旧样式 GB/T 7714-2015 (兼容保留)

%----- 参考文献引用格式 -----
\usepackage[numbers,sort&compress]{natbib}   % 顺序编码制引用
%\usepackage[numbers,super,square,sort&compress]{natbib}  % 上标引用
%\usepackage[authoryear,sort&compress]{natbib}            % 著者-出版年制
%\bibpunct{（}{）}{；}{a}{}{，}  % 著者-出版年制如需中文全角括号, 取消本行注释 (须配合 authoryear 引用格式)
\renewcommand{\bibfont}{\small}              % 参考文献正文字号
\setlength{\bibsep}{7pt plus 3pt minus 3pt}  % 参考文献条目间距

%----- 制作索引 -----
\usepackage{imakeidx}
\makeindex
\indexsetup{level=\chapter*}  % 索引标题的级别与章标题一致
% 提示: 编译环境若禁用了 shell 执行, makeindex 无法自动运行, 索引将为空
\ifnum\shellescape=0\relax
    \PackageWarningNoLine{zhbook}{编译时禁用了 shell 执行, 索引将不会自动生成;^^J请允许 shell 执行(TeX Live 默认即允许运行 makeindex), 或手动运行 makeindex main.idx}
\fi

% ===============================================================
%  四、按需微调 (取消注释后生效)
% ===============================================================
%\setlength{\itemsep}{3pt}           % 列表项之间的额外间距
%\raggedbottom                       % 允许页面底部参差, 避免出现过大的竖直空白
%\includepdf[pages=-]{pdfname.pdf}   % 插入外部 PDF 文件 (如自行制作的封面)
%\includeonly{chapters/chap03}           % 增量编译: 只编译改动过的章节 (首次需完整编译一次)
